\chapter{Teoretický úvod}
\label{chap:theory}

\section{Počítačem podporovaná teorie grafů}
\label{sec:computer-assisted-graph-theory}

Počítače se v teorii grafů nepoužívají pouze k numerickým výpočtům, ale také jako
nástroj pro objevování a ověřování matematických tvrzení. V počítačem podporované
teorii grafů můžeme například prohledávat grafy dané velikosti, vybírat z nich
grafy se zadanými vlastnostmi a testovat na nich navržená tvrzení.

Typickou úlohou je hledání protipříkladů k tvrzením tvaru: každý graf s vlastností
\(P_1\) má také vlastnost \(P_2\). Pokud takové tvrzení neplatí, chceme najít graf,
který splňuje \(P_1\), ale nesplňuje \(P_2\). Takový graf pak tvoří konkrétní
protipříklad k původnímu tvrzení.~\cite[Section~2]{jooken2025computerassisted}.

Takový postup sám o sobě nenahrazuje matematický důkaz obecného tvrzení. Je však
užitečný ve fázi hledání hypotéz a protipříkladů. Pokud nemůžeme tvrzení vyvrátit na
malých grafech, může nám počítač pomoci odhadnout jeho správnou formulaci a nasměrovat
hledání důkazu. Pokud naopak najdeme malý protipříklad, může nám pomoci pochopit, proč
tvrzení neplatí.


\section{Generování grafů}
\label{sec:graph-generation}

Generováním grafů budeme rozumět postup, při kterém algoritmus vytváří grafy
z předem určené třídy. V nejjednodušším případě bychom pro pevný počet vrcholů
mohli projít všechny možné hrany mezi těmito vrcholy. Pro graf na \(n\) vrcholech
existuje \(\binom{n}{2}\) možných hran, a tedy \(2^{\binom{n}{2}}\) označených
neorientovaných grafů. Toto číslo roste velmi rychle, takže přímé procházení všech
možností použijeme jen pro malé hodnoty \(n\).

V mnoha případech nám navíc nezáleží na konkrétních jménech vrcholů. Dva grafy,
které se liší pouze přejmenováním vrcholů, představují z hlediska mnoha vlastností grafů stejný 
případ.
Pokud bychom generovali všechny označené grafy, jeden graf bychom zpracovávali v mnoha
různých případech.
Praktické generátory se proto snaží vytvořit pouze jeden
reprezentant z každé izomorfní třídy.

Generátor obvykle nepoužívá jednoduchý seznam všech možností, ale generuje graf postupně.
Může například začít prázdným grafem a přidávat vrcholy nebo hrany a průběžně kontrolovat,
zda částečně vytvořený graf ještě může vést ke grafu požadované třídy.
Díky takovému inkrementálnímu  postupu můžeme zahazovat celé větve výpočtu dříve, než je dokončíme.

\section{Izomorfismus grafů}
\label{sec:graph-isomorphism}

Dva grafy \(G\) a \(H\) nazveme izomorfní, pokud existuje bijekce mezi jejich vrcholy,
která zachovává sousednost. Jinými slovy, grafy mají stejnou strukturu a liší se pouze
pojmenováním vrcholů. Při generování grafů je izomorfismus zásadní,
protože nechceme opakovaně zpracovávat různé popisy téhož grafu.

Pokud grafy obsahují další strukturu, například barvy vrcholů nebo vybrané zakotvené
vrcholy, musí izomorfismus respektovat i tyto informace. U obarvených grafů tedy můžeme
zaměňovat pouze vrcholy se stejnou barvou. Tato myšlenka je pro nás důležitá, protože v této práci rozšiřujeme generátor
\texttt{geng} o barevné třídy a zakotvené vrcholy, které musí izomorfismus
v procesu generování respektovat.

Prakticky narazíme na to, že testování izomorfismu je výpočetně náročná úloha
a při generování se může objevit velmi často. Efektivní generátory proto používají
specializované algoritmy a reprezentace, které umožňují rychle rozhodovat, zda jsme daný
případ již nevygenerovali v jiné podobě.

\section{Kanonické reprezentace}
\label{sec:canonical-representations}

Jedním ze způsobů, jak pracovat s izomorfismem, je přiřadit každému grafu
kanonickou reprezentaci. Chceme, aby všechny izomorfní kopie téhož grafu dostaly stejný
výsledek. Pokud tedy dva grafy mají stejnou kanonickou reprezentaci, můžeme je považovat za navzájem izomorfní.

Kanonické značení se používá například při ukládání výsledků, porovnávání grafů
nebo odstraňování duplicit. Generátory ho využívají uvnitř samotného algoritmu,
aby nevytvářely více zástupců jedné izomorfní třídy. Jedním z takových nástrojů pro
práci s kanonickým značením grafů je knihovna nauty.~\cite{mckay2014practical}.

Kanonické značení se komplikuje ve chvíli, kdy graf nenese pouze informaci o hranách, ale
také další data. Pokud mají vrcholy barvy, musí kanonické značení zachovat rozdělení vrcholů
do barevných tříd. Při návrhu rozšíření pro obarvené grafy proto nestačí přidat barvy až do výstupu.
Musíme je zahrnout už do testů izomorfismu.

\section{Backtracking a prořezávání}
\label{sec:backtracking-pruning}

Mnoho generátorů grafů je založeno na backtrackingu. Algoritmus postupně vytváří částečný graf
a z něj zkouší všechny přípustné možnosti rozšíření. Pokud zjistíme, že daný částečný graf nemůže
vést k žádnému požadovanému výsledku, vrátíme se a pokračujeme jinou větví výpočtu.

Tomuto předčasnému zahazování větví říkáme prořezávání. Prořezávání může být založeno na
jednoduchých podmínkách, například na počtu hran nebo stupních vrcholů, ale také na složitějších
algoritmech a heuristikách. Čím dříve dokážeme větev zahodit, tím méně práce musí generátor vykonat.

V kontextu této práce budeme rozlišovat filtrování hotových grafů a prořezávání částečně
vygenerovaných grafů. Filtrování hotových grafů je jednodušší, protože pracujeme s kompletním
grafem. Prořezávání během generování může být výrazně efektivnější, ale musíme počítat s tím,
že pracujeme pouze s částečným grafem a některé vlastnosti ještě nemusí být rozhodnutelné.

\section{Praktické komplikace generování grafů}
\label{sec:practical-complications}

První praktickou komplikací je kombinatorická exploze. I když generujeme pouze neizomorfní grafy,
počet výsledných grafů může být velmi vysoký. Generátory proto musí být implementovány s důrazem na efektivitu
a často používají nízkoúrovňové datové struktury, například bitové reprezentace množin hran.

Druhou komplikací je práce s interní reprezentací grafu. Výkonné generátory jsou často psané
v jazyce C a jejich vnitřní struktury jsou optimalizované pro rychlost, nikoli pro pohodlné
použití z aplikačního kódu. Pokud chceme nad generovanými grafy spouštět vlastní algoritmy,
často musíme psát pluginy nebo pracovat přímo s detaily implementace generátoru.

Třetí komplikací je propojení generování s dalšími algoritmy. V experimentální práci
běžně chceme každý vygenerovaný graf ihned testovat, obarvovat, porovnávat nebo převádět do jiného 
formátu. Pokud kvůli tomu musíme graf kopírovat do jiné datové struktury, můžeme ztratit podstatnou
část výkonu původního generátoru.

Tyto problémy motivují návrh rozhraní, které zachová výkon specializovaných generátorů, ale 
současně poskytne pohodlnější způsob konfigurace, filtrování a zpracování grafů.
